\documentclass[11pt,a4paper]{article}
\usepackage[margin=24mm]{geometry}
\usepackage{amsmath,amssymb,amsthm,booktabs,longtable,hyperref}
\setlength{\emergencystretch}{2em}
\newtheorem{theorem}{Theorem}\newtheorem{lemma}[theorem]{Lemma}\newtheorem{proposition}[theorem]{Proposition}\newtheorem{corollary}[theorem]{Corollary}
\newcommand{\Aut}{\operatorname{Aut}}\newcommand{\PSL}{\operatorname{PSL}}\newcommand{\PSU}{\operatorname{PSU}}\newcommand{\PSp}{\operatorname{PSp}}\newcommand{\Om}{\Omega}\newcommand{\Sp}{\operatorname{Sp}}
\newcommand{\ii}{i_2}\newcommand{\nuu}{\nu_2}
\title{Chiral generating pairs in finite simple groups:\\a complete assembly of the positive theorem}
\author{Human Compression project}\date{}
\begin{document}\maketitle
\begin{abstract}
We assemble a proof of the positive existence theorem of Leemans and
Liebeck for finite nonabelian simple groups. The proof combines explicit
alternating generators, finite word certificates, and involution-counting
arguments with all parameter ranges accounted for. The classification
and geometric reconstruction theorems used as antecedents are stated
explicitly. Several auxiliary assertions in the original proof are replaced
by corrected arguments, without changing its positive conclusion. No new
existence theorem or classification is claimed.
\end{abstract}
The full proofs and certificates used below are retained in their exact
component editions. The accepted general orthogonal counting argument
and the separately accepted $D_4^+$ cover are joined explicitly in
Section~4.4. Historical pending-review language in a frozen component
records its earlier author stage; the acceptance and dependency records
bound to this assembly state the completed interfaces. This manuscript
establishes the positive principal theorem, with the source boundaries
specified in Section~5.

\section{The exact theorem and its geometric meaning}
The fixed source is arXiv:1603.07713v3, whose running header is dated
18 May 2016 and whose internal date is 6 November 2018. Its numbered
principal theorem is the following positive statement~\cite{LL}.
\begin{theorem}[Leemans--Liebeck]\label{main}
Let $G$ be a finite nonabelian simple group which is not isomorphic to
$A_7$, $\PSL_2(q)$, $\PSL_3(q)$ or $\PSU_3(q)$.
There are $x,t\in G$ such that
\[
 G=\langle x,t\rangle,\qquad |t|=2,
 \qquad\nexists\alpha\in\Aut(G):\quad
 x^\alpha=x^{-1},\quad t^\alpha=t.
\]
\end{theorem}
The source asks only that no \emph{involutory} automorphism have the
specified action. These formulations are equivalent for a generating
pair: the square of such an automorphism fixes both generators and hence
is the identity. The identity itself cannot invert $x$, since $|x|>2$.
Indeed, a group generated by two involutions is dihedral, whereas $G$ is
nonabelian simple. The automorphism group in the theorem is the full
$\Aut(G)$; restricting to inner automorphisms would not suffice.

\begin{lemma}\label{geometry}
Every pair in Theorem~\ref{main} satisfies the rank-three rotation
intersection conditions and yields an abstract chiral polyhedron with
full automorphism group $G$.
\end{lemma}
\begin{proof}
Set $y=x^{-1}t$. Then $xy=t$, so $(xy)^2=1$, and
$G=\langle x,y\rangle$. An element of
$\langle x\rangle\cap\langle y\rangle$ commutes with both generators,
and therefore is central in $G$. This intersection is trivial.
Neither $x$ nor $y$ has order at most two: otherwise $G$ would be cyclic
or generated by two involutions. Thus the orders are at least three.

We now invoke the retained Schulte--Weiss rank-three reconstruction
and regularity criterion~\cite{SW}, in the form used in the source:
these rotation data reconstruct a directly regular or chiral polyhedron;
the directly regular alternative is precisely the existence of the
reversing automorphism in Theorem~\ref{main}. Its absence gives chirality,
and the reconstructed full automorphism group is $G$.
\end{proof}
\begin{corollary}\label{polyhedra}
Every group in Theorem~\ref{main} is the full automorphism group of an
abstract chiral polyhedron.
\end{corollary}
This uses geometric reconstruction as an explicit antecedent, rather
than claiming a new construction theorem for abstract polytopes.
Equivalence of marked rotation data means a group isomorphism carrying
the distinguished generators to their counterparts. No bijection between
\emph{literal labelled pairs} and unmarked polyhedra is asserted here.
In particular, no enumerative count is deduced by forgetting orientation,
identifying enantiomorphs, or dividing the number of pairs by an unstated
symmetry factor.

\section{Two sufficient criteria}
The following elementary arguments explain the two forms of certificate
used throughout the proof. Write $\ii(H)$ for the number of elements of
order exactly two, and $\nuu(H)=\ii(H)+1$ for the number of solutions of
$h^2=1$.
\begin{lemma}[Word criterion]\label{word}
If $G=\langle x,t\rangle$, $|t|=2$, and a word $w$ satisfies
$|w(x,t)|\ne|w(x^{-1},t)|$, no automorphism of $G$ inverts $x$ and fixes $t$.
\end{lemma}
\begin{proof}
An automorphism with those generator images sends one word value to the
other, whereas automorphisms preserve element order.
\end{proof}
All forty-seven explicit finite certificates used below employ
$w(x,t)=xtx^2tx^3t$: twelve sporadic cases and thirty-five Lie parameter
entries. These are parameter certificates, not forty-seven distinct
abstract groups, since some small classical names coincide.

\begin{lemma}[Counting criterion]\label{count}
Let $G$ be finite and centreless, and fix $x\in G$ of order greater than two.
Suppose $\mathcal H$ is a family of actual proper subgroups covering every
proper subgroup containing $x$. If
\[
 \ii(G)>
 \sum_{\substack{\alpha\in\Aut(G)\\|\alpha|=2,\ x^\alpha=x^{-1}}}
       \ii(C_G(\alpha))+
 \sum_{H\in\mathcal H}\ii(H),
\]
then an involution $t$ satisfies the conclusion of Theorem~\ref{main}.
\end{lemma}
\begin{proof}
The second sum bounds involutions that generate a proper subgroup with
$x$. The first bounds involutions fixed by any involutory reverser.
Choose $t$ outside their union. Any reverser fixing this generating
partner has square one, as observed above, and is already excluded.
\end{proof}
Overlaps among these sets only enlarge the bound, so no disjointness is
required. A list of \emph{types} or conjugacy classes of subgroups is not
yet $\mathcal H$: the number of actual subgroups containing $x$ must be
proved. The cyclic normalizer must also be covered or separately charged.
These points are paid explicitly by the acquired components.

A convenient bound for the first sum is the following. If
$G\le J\lhd\Aut(G)$ and $C_J(x)=T$, each outer coset contains at most
$|T|$ reversers. Only square-one elements of $\Aut(G)/J$ can be images
of involutory reversers. If their number is $\kappa$ and every relevant
fixed group has at most $\mu$ involutions, the first sum is at most
$\kappa|T|\mu$. This does not assume that all reversers are conjugate.
When passing through a kernel $K$, we use
$\nuu(H)\le |K|\nuu(H/K)$, including the identity before taking fibres.
No algebraic dimension estimate is substituted for a finite kernel order.

\section{Alternating and sporadic groups}
For $A_n$, $n\ge8$, the complete alternating component~\cite{Alt}
proves generation and excludes the full reversing automorphism group
for the source's explicit pairs:
\[
\begin{array}{c|cc}
n\text{ even}&x=(2,3,\ldots,n)&t=(1,2)(3,4),\\
n\text{ odd}&x=(1,2,\ldots,n)&t=(1,2)(3,6).
\end{array}
\]
Its proof is uniform in $n$; the finite controls only check displayed
identities and do not prove an unbounded theorem by sampling. The
standard alternating-group automorphism theorem and the stated primitive
permutation-group antecedent remain retained. The lower alternating
names $A_5\cong\PSL_2(4)$ and $A_6\cong\PSL_2(9)$ lie in the excluded
isomorphism families, as does $A_7$ by name.

The twenty-six sporadic groups are divided into the disjoint sets
\[
\begin{split}
\mathcal S_{12}={}&\{M_{11},M_{12},M_{22},M_{23},M_{24},J_1,J_2,J_3,
 HS,He,Suz,McL\},\\
\mathcal S_{14}={}&\{J_4,Co_1,Co_2,Co_3,Fi_{22},Fi_{23},Fi'_{24},
 M,B,Th,Ly,\mathrm{O'N},HN,Ru\}.
\end{split}
\]
For $\mathcal S_{12}$, the exact permutation certificates~\cite{Twelve}
check membership in each identified named group, full generation, and
the unequal word orders of Lemma~\ref{word}. The $M_{23},M_{24}$ source
permutations are retained literally. Other witnesses may have different
polyhedron types from the source examples; their role is existence.
Named representation identifications and simplicity are retained inputs,
not inferred from an order match alone.

For $\mathcal S_{14}$, the complete counting component~\cite{Sporadic}
chooses $x$ of an indicated odd prime order $p$. Counting conjugates of a
maximal subgroup $M$ containing $x$ gives the bound
\[
 B_G(x)=|C_G(x)|\sum_{[M]}\frac{i_p(M)\ii(M)}{|M|}.
\]
The sum includes every maximal-subgroup conjugacy class, including
repeated character-table names for distinct classes. The inverter term
is zero if $x$ and $x^{-1}$ are not conjugate in the full automorphism
group; otherwise it is at most $|C_{\Aut(G)}(x)|\mu$.
The component supplies a rigorous $\mu$ from complete maximal lists in
$G$ or its full automorphism group, checking the required centralizer
indices in the four outer-index-two cases. Its fourteen strict rational
inequalities prove Lemma~\ref{count}. Ordinary character tables, class
fusions, maximal-subgroup completeness, and the named automorphism-group
identifications are explicit retained classification premises. Thus all
twenty-six sporadic cases are covered without using any historical
whole-map census or its numerical total.

\section{The Lie families and all exceptional parameters}
We use the standard classification of finite nonabelian simple groups,
including its low-rank isomorphisms, to pass from these families to all
simple groups. In the classical notation below, groups are projective
where the centre is nontrivial. Orthogonal odd dimension has odd $q$;
in characteristic two the corresponding simple groups are symplectic.

\subsection{Finite parameter certificates}
The finite component~\cite{Finite} supplies every simple entry in the
source's two finite exception lists, together with its three exceptional
Lie cases. For $\epsilon=\pm$, $L^+=\PSL$ and $L^-=\PSU$, these are
\[
\begin{gathered}
 L_4^\epsilon(2),L_4^\epsilon(3),L_4^\epsilon(4),
 L_5^\epsilon(2),L_6^\epsilon(2),\quad
 \PSp_4(3),\Sp_4(4),\Sp_6(2),\PSp_6(3),\Sp_8(2),\\
 \Om_7(3),\Om_8^\epsilon(2),P\Om_8^+(3),\Om_{10}^\epsilon(2),\\
 L_5^\epsilon(3),L_6^\epsilon(3),L_8^\epsilon(2),
 \Sp_6(4),\PSp_8(3),\Sp_8(4),\Om_{12}^\epsilon(2),\\
 G_2(3),G_2(4),{}^2F_4(2)'.
\end{gathered}
\]
There are $21+11+3=35$ parameter entries. The source also prints
$\Sp_4(2)\cong S_6$, which is not simple and is not a positive obligation.
The overlaps $\PSL_4(2)\cong A_8$ and
$\PSU_4(2)\cong\PSp_4(3)$ are preserved as duplicate parameter routes,
not counted as new abstract groups. Each positive certificate verifies
its actual named-constructor interface, membership, full generation and
Lemma~\ref{word}; a candidate-producing search is not the correctness
premise.

\subsection{Exceptional Lie type}
The complete exceptional component~\cite{Exceptional} treats
\[
 E_8,E_7,E_6,{}^2E_6,F_4,G_2,{}^3D_4,{}^2F_4,{}^2G_2,{}^2B_2
\]
in every simple parameter not covered by the finite list. The small
identifications $G_2(2)'\cong\PSU_3(3)$ and
${}^2G_2(3)'\cong\PSL_2(8)$ belong to excluded families and are not
converted into positive cases. The Tits group is supplied by the finite
certificate, not by treating ${}^2F_4(2)$ as simple.

The component proves regularity of the chosen torus elements by a
root-lattice argument, including its small finite cases. It bounds
reversers by outer cosets, identifies the admissible algebraic
involution classes using the reviewed flag-variety argument, and charges
the cyclic normalizer. Explicit class orders give the involution lower
bounds; translated polynomials with nonnegative coefficients and positive
constant terms prove all uniform margins. The residual fixed parameters
are checked exactly. Torus existence, complete maximal-overgroup types,
class centralizers and the uniform Lawther--Liebeck--Seitz upper theorem~\cite{LLS}
remain named antecedents. In particular, the proof does not invoke
Seitz's theorem outside its field-size hypotheses or infer a finite
lower constant from an asymptotic estimate.

\subsection{Linear, unitary and symplectic type}
The elementary classical lower component~\cite{Lower} proves
\[
\begin{array}{c|c}
G&\ii(G)>q^D/4\\\hline
\PSL_n(q),\PSU_n(q),\ n\ge4&D=\lfloor n^2/2\rfloor\\
\PSp_{2m}(q),\ m\ge2&D=m^2+m\\
\Om_{2m+1}(q),\ m\ge3&D=m^2+m\\
P\Om_{2m}^{\epsilon}(q),\ m\ge4&D=m^2-1.
\end{array}
\]
It counts actual involution families with their scalar losses; no
maximal-subgroup classification is needed for these lower estimates.
For linear and unitary groups, the complete upper component~\cite{Linear}
constructs the Singer or $(n-1)+1$ element, proves its adjoint torus
centralizer and the single possible reversing coset, and bounds the
actual fixed groups, including characteristic-two nonalternating forms.
It then proves the strict counting inequality for every parameter outside
the displayed finite list. The exact source Table 4 linear/unitary cover
is explicitly retained as a named classification assertion, with its
Bereczky--Bamberg--Penttila provenance and the published author erratum.
The later acquisition of an older full BP preprint does not retroactively
turn that frozen component into an independent proof of the classification.

For symplectic groups, the complete component~\cite{Symplectic} uses the
full Singer cover in Bubboloni--Spiga--Weigel Lemma 5.1~\cite{BSW}, including the
odd-characteristic unitary overgroup omitted from the source table.
Trace forms classify the natural reversers; the exceptional graph-field
coset in degree four and even characteristic is separately included.
Actual finite kernels and the degree-four estimates give strict
inequalities in all remaining ranks and fields. The finite list supplies
all small excluded parameters. Thus both components meet exactly the
hypotheses of Lemma~\ref{count}.

\subsection{Orthogonal type}
The torus component~\cite{Tori} supplies elements with their actual
spinor, Dickson and projective scalar conditions. In odd dimension use
a squared primitive norm-one multiplier on a negative hyperplane; in
minus even dimension use a primitive multiplier for even $q$ and its
square for odd $q$; in plus dimension use primitive multipliers on the
negative $(2m-2)+2$ decomposition. The connected projective conformal
centralizer orders are respectively
\[
 q^m+1,\qquad q^m+1,\qquad(q^{m-1}+1)(q+1).
\]
The odd and even fixed-group components~\cite{Odd,Even} pay the finite
centralizers of natural reversing automorphisms, including nonsquare
multipliers and unipotent kernels. In particular the full-rank even
characteristic degree-eight bound is $\nuu(C_G(\alpha))\le q^{11}$.

The orthogonal counting component~\cite{Orthogonal} uses the complete
BSW overgroup statements in odd dimension, minus dimension, and plus
dimension at least ten. It proves the multiplicities of actual field
structures, includes a bound for the cyclic normalizer, and establishes
the uniform inequalities. Its remaining parameters are exactly those
already present in the finite list, except for the separately treated
dimension-eight range.

In plus dimension eight, the complete D4 specialization~\cite{Dfour}
proves, for every primitive choice on the negative $6+2$ decomposition,
that the only maximal types are its decomposition stabilizer, at most
two unitary extension normalizers, and the triality $B_3$ subgroups.
It derives the cover from the complete published BHRD Table 8.50~\cite{BHRD};
the table's number of conjugacy classes is not used as an element-containing
multiplicity. The independent triality argument in~\cite{Orthogonal}
chooses the primitive components so that all $B_3$ overgroups are absent
for even $q\ge4$ and odd $q\ge7$. At $q=5$ it counts the actual six
possible line stabilizers, not the source table's asserted two.
The same argument controls the additional inverter cosets.
With the exact degree-eight lower estimates, the displayed strict
inequalities prove the claim at all these fields. The finite certificates
pay plus type at $q=2,3$ and minus type at $q=2$.
This is the precise interface at which the complete D4 cover discharges
the conditional hypothesis of the counting manuscript; no later result
is used to justify that cover.

\section{Completion and source boundaries}
\begin{proof}[Proof of Theorem~\ref{main}]
By the retained classification of finite simple groups, every group under
the hypothesis is alternating, sporadic or of Lie type, with the standard
small identifications described above. The alternating proof covers
all eligible alternating groups. The two disjoint sporadic rosters cover
all twenty-six sporadics. The finite list, exceptional component and
three classical routes cover every eligible Lie parameter. Each route
establishes the same full-automorphism generating-pair conditions, either
by Lemma~\ref{word} or Lemma~\ref{count}. There is no assumption that
an auxiliary table is itself a theorem about an unverified new element.
This proves the group-theoretic assertion; Lemma~\ref{geometry} proves
Corollary~\ref{polyhedra}.
\end{proof}

The manuscript establishes this \emph{positive} theorem. On source page 3,
nonexistence for $\PSL_2(q)$ is attributed to Macbeath and Singerman,
and nonexistence for $A_7$ to the earlier atlas computation. These are
historical antecedents, not deductions of our positive construction.
The corresponding $\PSL_3(q),\PSU_3(q)$ statement is explicitly announced
there as forthcoming work [19], not proved in the fixed source. We do
not promote that announcement into a proved converse. The source's
introductory classification of regular polyhedron groups, and its final
almost-simple Suzuki and $\PSL_2$-socle statements, likewise remain
attributed antecedents. The general almost-simple extension is a question,
not an additional conclusion.

The Section 6 counts for $M_{11},M_{12},M_{22},J_1,J_2$, the cited
all-types O'Nan enumeration, and example types not realized by our
replacement witnesses are auxiliary historical statements. Our theorem
requires their existence conclusion only. No new exhaustive map count,
faithfulness of an unrelated historical census, or claim to recover every
ancillary numerical example is made. In particular the O'Nan existence
argument does not import the disputed total from the separate L046 source.

Several source proof interfaces have been replaced: the omitted normalizer
term is covered; outer inverters are not assumed to have one class; small
field hypotheses are respected; the symplectic unitary overgroup is included;
orthogonal scalar losses and triality multiplicities are corrected.
The source's special estimate $\ii(\PSU_4(q))>q^8/2$ is false for
$q\equiv3\pmod4$, as established by the exact formula in~\cite{Lower}.
The valid quarter bound and sharper finite calculations suffice here.
These are repairs to supporting assertions; they do not refute the
positive principal theorem.

\paragraph{Attribution and retained costs.}
The theorem, principal strategy and alternating pairs belong to
Leemans--Liebeck and their cited predecessors. The full component proofs
and exact certificates listed below are part of this assembly, not merely
status summaries. The classification of finite simple groups, Schulte--Weiss
reconstruction, named finite-group constructions and character tables,
complete maximal-subgroup classifications, stated automorphism structure,
and the cited Lie-type involution inputs remain explicit borrowed results.
They are not claimed to have been reproved or formally verified.

The root agent authored the alternating, twelve-sporadic, sporadic counting, exceptional,
orthogonal torus, odd fixed-group, triality and orthogonal counting work.
Publication\_recovery authored the finite Lie certificates, classical
lower bounds, linear/unitary, symplectic, even fixed-group and D4 cover
components and this assembly. The root supplied the finite candidates and
several organizing observations; l044\_resume supplied the exceptional
flag-variety amendment, which has its own separate review. The two
relative-field arguments were developed in parallel and are not counted
as independent reviews of each other. f003\_resume supplied the bounded
order-$21/63$ note and distinctly reviewed the general counting/multiplicity component.
Source-contract assistance by l044\_resume is recorded separately from
mathematical acceptance. Earlier bounded proofs, their distinct reviews,
and all retained classification costs remain attached to their exact
editions. No novelty, new question resolution, Lean verification or new
measurement is asserted.

\begin{thebibliography}{99}\small
\bibitem{LL} D. Leemans and M. W. Liebeck, \emph{Chiral polyhedra and finite
simple groups}, arXiv:1603.07713v3, fixed fourteen-page source.
\bibitem{SW} E. Schulte and A. Ivi\'c Weiss, \emph{Chiral polytopes},
DIMACS Series 4 (1991), 493--516; retained rank-three reconstruction.
\bibitem{Alt} \emph{Alternating groups generated by a cycle and two transpositions}, complete acquired
L045 alternating component, version 2, with distinct review and root adoption.
\bibitem{Twelve} \emph{Chirality detected by the order of one word}, complete
acquired twelve-sporadic component, version 2, with distinct review and adoption.
\bibitem{Sporadic} \emph{Counting involutions that generate a chiral polyhedron},
complete acquired fourteen-sporadic component, version 1, with distinct review.
\bibitem{Finite} \emph{Finite Lie-group witnesses for chiral polyhedra}, complete
acquired thirty-five-parameter component, version 1, with distinct review.
\bibitem{Exceptional} \emph{Involutions outside the reversal and subgroup
obstructions}, exceptional component, version 2, including the separately
reviewed flag-variety amendment and complete arithmetic certificates.
\bibitem{Lower} \emph{Explicit involution classes for the classical part of L045},
version 1, with exact unitary correction and distinct review.
\bibitem{Linear} \emph{Torus reversers and explicit bounds for the linear and
unitary branches}, version 1, with distinct review.
\bibitem{Symplectic} \emph{Explicit symplectic reverser and overgroup bounds},
version 2, with distinct review.
\bibitem{Tori} \emph{Orthogonal tori and their projective orders}, version 1, with distinct review.
\bibitem{Odd} \emph{Bounds for orthogonal fixed groups in odd characteristic},
version 2, with distinct review.
\bibitem{Even} \emph{Finite kernels of full-rank orthogonal inverters in characteristic two},
version 1, with distinct root review.
\bibitem{Dfour} \emph{Maximal overgroups of a full negative $6+2$ torus element},
version 1; complete type specialization distinctly reviewed, shared multiplicity
interface separately routed through the counting review.
\bibitem{Orthogonal} \emph{Counting overgroups of an orthogonal torus element}
and \emph{Choosing an orthogonal torus element to separate the triality modules}, complete orthogonal counting component, version 1;
with its separate distinct review and the independently accepted complete $D_4^+$ cover.
\bibitem{LLS} R. Lawther, M. W. Liebeck and G. M. Seitz,
\emph{Fixed point ratios in actions of finite exceptional groups of Lie type},
Pacific J. Math. 205 (2002), 393--464, especially Proposition 1.3.
\bibitem{BSW} D. Bubboloni, P. Spiga and Th. Weigel,
\emph{Normal $2$-coverings of the finite simple groups and their generalizations},
arXiv:2208.08756v1, Lemmas 5.1, 6.1, 7.3 and 8.1.
\bibitem{BHRD} J. N. Bray, D. F. Holt and C. M. Roney-Dougal,
\emph{The Maximal Subgroups of the Low-Dimensional Finite Classical Groups},
LMS Lecture Note Series 407, Cambridge University Press, 2013,
Main Theorem 2.1.1 and Table 8.50.
\end{thebibliography}
\end{document}
